\section{Evaluation Methodology}
\label{sec:method}

\subsection{Dataset and the deployed model's provenance}
We evaluate detection on N-BaIoT~\cite{meidan2018nbaiot}: benign traffic from
nine commercial devices (doorbells, a thermostat, a baby monitor, security
cameras and a webcam) and traffic from the same devices after infection by
Mirai and BASHLITE, each executing five attacks. Features are the 115 Kitsune
statistics described in \S\ref{sec:design}. The attack classes are defined by
the dataset's authors; for example, BASHLITE \emph{COMBO} is ``sending spam
data and opening a connection to a specified IP address and port'' and Mirai
\emph{UDPplain} is ``UDP flooding with fewer options, optimized for higher
PPS''~\cite{meidan2018nbaiot}.

The model shipped with \textsc{Sentinel} had been trained in a notebook that
was not kept, so its training data was undocumented. We recovered it from the
saved min--max scaler: its 115 per-feature minima and maxima match the benign
traffic of exactly six of the nine devices with zero deviation, and the number
of rows it was fitted on (\NprovRows{}) equals their combined benign row count.
We call these six devices \emph{in-distribution} and the remaining three
(\emph{Ennio} doorbell, \emph{Philips} baby monitor, \emph{SimpleHome} XCS7-1003
camera) \emph{held-out}. No model evaluated here saw any traffic from the
held-out devices, which therefore measure how the detector behaves on a device
it has never observed---the ordinary case for a home appliance.

\subsection{Splits and leakage}
Adjacent rows of Kitsune features share damped statistics, so a random
train/test split places near-duplicate rows on both sides of the boundary.
Following Meidan~et~al., we split each in-distribution device's benign
traffic \emph{chronologically} and keep the final third for
testing~\cite{meidan2018nbaiot}. We train on the first 90\% of the remainder
and use the last 10\% for early stopping and threshold selection. Scalers are
fitted on the training portion only. Attack traffic is never used for training
or threshold selection. Windows are ten consecutive rows from a single capture
file, never spanning a file or device boundary; test windows do not overlap.

The deployed model does not satisfy these conditions---its scaler saw all
benign rows of the in-distribution devices---so we report its in-distribution
false-positive rate only as optimistic and rely on its held-out-device and
attack results, which are clean.

\subsection{Models compared}
\begin{itemize}
\item \emph{Deployed}: the shipped ONNX model with its shipped threshold.
\item \emph{LSTM-AE}: the same architecture retrained under the protocol above
  with five seeds. Its parameter count reproduces the deployed model's exactly.
\item \emph{Dense AE}: the per-row deep autoencoder architecture of N-BaIoT's
  authors~\cite{meidan2018nbaiot} (hidden widths 86--57--37--28, mirrored),
  scoring a window by its rows' mean error. It sees the same windows as the
  LSTM-AE but models no temporal order.
\item \emph{Isolation Forest}~\cite{liu2008iforest} on the flattened window,
  as a non-neural baseline.
\end{itemize}
Every retrained model uses the same threshold rule: the 99th percentile of its
own validation-benign scores, a nominal 1\% false-positive budget fixed without
reference to attack traffic.

\subsection{LLM classification (pre-registered)}
To avoid grading our own test cases after seeing the model's answers, we
committed the experimental design to the project repository before running the
model on any sample. The ground truth maps each N-BaIoT attack class, by its
authors' definition, to a set of acceptable ATT\&CK techniques: scans to
Network Service Discovery, Active Scanning or Remote System Discovery, and all
floods to Network or Endpoint Denial of Service. Equivalent techniques from the
ATT\&CK for ICS matrix are accepted, because \textsc{Sentinel}'s prompt asks
for that matrix. We drew 20 windows per attack class (200 in total) with a fixed
seed, from both device groups, and queried the model under two prompts:
\emph{as-built}, produced by \textsc{Sentinel}'s own prompt builder, and
\emph{evidence}, identical but for five plain-language traffic statistics
computed from the window (recent packet count, mean packet size, approximate
numbers of destination hosts and of sockets, and channel jitter). Fields that
could leak the label---the protocol, notably---were withheld from both, and
reputation lookups were disabled. We report technique and tactic accuracy with
Wilson 95\% intervals~\cite{wilson1927}, the rate of valid JSON, and latency.
Two further analyses were declared after the initial design but before any
output was seen: as a secondary analysis, results restricted to the windows the
detector would actually forward to the classifier; and, labelled exploratory,
the rate at which the pipeline would have blocked, i.e.\ the fraction of
attack windows assigned severity~$\geq 6$.

\subsection{Latency and host}
We time each stage of the path in process with the real components: feature
extraction per packet on real traffic (the Mirai capture distributed with
Kitsune~\cite{kitsunepy}), anomaly inference per window, a Redis bus hop, the
response policy, and the threat-intelligence lookups. LLM latency is taken from
the \NllmLatN{} classification calls that ran on this host. All timing
measurements ran on one host: a four-core
aarch64 virtual machine with 7.7\,GB of memory. This is the same instruction
set as, and a similar core count and memory size to, the Raspberry~Pi~5 that
\textsc{Sentinel} targets, but it is not that device: it does not reproduce the
Pi's clock speed, memory bandwidth or thermal throttling, so absolute timings
on the Pi will differ. Neural models were trained on a Colab T4 GPU.

\emph{Deviation from the pre-registration.} The first \NllmLocN{} of the
400 LLM calls ran on this host as planned; the run was then stopped at a
checkpoint because it monopolised the workstation, and the remaining
\NllmCloudN{} ran on a cloud GPU with the same model (verified by digest),
prompts and sampling. We recorded this, and its consequences, in the
repository before any cloud output existed: latency is reported from the
local calls only, and accuracy pools both hosts, with a per-host check.
